\chapter{Computing systems}
\label{chapter-computing-systems}

\section{What I am trying to learn}
\label{section-computing-systems-goal}

The immediate goal is not to become a
computer engineer.  It is to understand
enough about computers to make good
choices about my own machines.

In particular, when buying a computer I
want to be able to answer:

\begin{itemize}
\item What work is limited by the CPU?
\item How much RAM do I actually need?
\item When does a GPU matter?
\item What is VRAM, and why does it matter
for local AI?
\item Which parts of the machine can I
replace or upgrade?
\item Which restrictions come from the
hardware, and which come from the
operating system?
\item What does it mean for a computer to
be open?
\item How much of my work can remain local
and private?
\end{itemize}

For my present situation there is one more
constraint: I move around a lot.  So the
main machine has to be a laptop.  A large
desktop workstation may make sense later,
but it is not the present problem.

\section{The basic hardware}
\label{section-computing-systems-hardware}

A computer has several components doing
very different jobs.

\medskip\noindent
{\bf CPU.}
The central processing unit executes most
ordinary programs.  Compiling code,
running a computer algebra system, and
many mathematical computations are
primarily CPU jobs.

A CPU has several cores.  More cores help
when a program can divide its work into
parallel tasks.  A fast individual core is
still important because many programs
cannot parallelize everything.

\medskip\noindent
{\bf RAM.}
RAM is the computer's short-term working
memory.  Programs and data being actively
used live there.

Running out of RAM does not merely mean
that no more files can be stored.  The
operating system starts moving data
between RAM and the much slower SSD.  This
can make a machine feel dramatically
slower.

For my work, RAM is particularly important
because I may simultaneously have a
browser, editor, TeX compiler, terminals,
computer algebra systems, and AI tools
open.

\medskip\noindent
{\bf Storage and SSD.}
The SSD is long-term storage.  It contains
the operating system, programs, source
files, papers, datasets, and models.

An SSD is much slower than RAM but much
faster than an old mechanical hard drive.
Capacity and replaceability are separate
questions: a large soldered SSD can still
be less useful in the long run than a
standard replaceable one.

\medskip\noindent
{\bf GPU.}
The graphics processing unit was designed
for massively parallel graphics
computations.  The same kind of
parallelism is extremely useful in machine
learning.

For ordinary LaTeX, browsing, Git, and
most programming, a powerful GPU is not
very important.  For training or running
many AI models it can become one of the
most important components.

\medskip\noindent
{\bf VRAM.}
A discrete GPU normally has its own memory,
called VRAM.  A model being run on the GPU
has to fit, at least to a large extent,
inside this memory.

Thus, for local AI, it is not enough to
ask whether a machine has a fast GPU.  The
amount of memory available to the GPU can
be the decisive limitation.

Some architectures instead use a common
pool of memory shared by the CPU and GPU.
This is called unified memory.  It changes
the comparison because a model may have
access to a much larger memory pool than
on a laptop GPU with a small fixed amount
of VRAM.

\medskip\noindent
{\bf Cooling and power.}
A fast chip produces heat.  A computer can
only sustain high performance if it can
remove that heat and supply enough power.

This is one reason laptop performance
cannot be judged from the processor name
alone.  The same chip can behave
differently in two machines with different
cooling systems and power limits.

\section{The software stack}
\label{section-computing-systems-software}

The hardware is only the bottom layer.

Very roughly, the stack is

$$
\begin{aligned}
&\text{applications}\\
&\text{operating system}\\
&\text{drivers and kernel}\\
&\text{firmware}\\
&\text{hardware}.
\end{aligned}
$$

The operating system manages files,
memory, processes, devices, networking,
permissions, and the interface between
applications and hardware.

The kernel is the central part of the
operating system.  Drivers teach the
kernel how to communicate with particular
hardware.

Firmware is software stored very close to
the hardware itself.  For example, a
laptop contains firmware controlling parts
of the boot process and many devices.

This distinction matters because
installing Linux changes a large part
of the software stack, but it does not
change the physical machine.  Linux cannot
make soldered RAM replaceable.

\section{What openness means}
\label{section-computing-systems-openness}

I should avoid using the word open as if
it meant only one thing.  There are
several different freedoms.

\begin{itemize}
\item {\bf Software freedom:} can I run,
inspect, modify, and replace the software?
\item {\bf Operating-system freedom:} can
I install another operating system?
\item {\bf Hardware freedom:} can I open
the machine and replace components?
\item {\bf Repair freedom:} can I obtain
parts and repair the machine without the
manufacturer's permission?
\item {\bf Data freedom:} are my files in
portable formats that I can move elsewhere?
\item {\bf Service independence:} does my
workflow continue to work if I stop using
a particular company's cloud services?
\end{itemize}

These freedoms are related but not
identical.  A machine can run open-source
software while having soldered hardware.
A repairable laptop can still run
proprietary software.  A Linux computer
can still depend heavily on cloud services.

\section{Local AI and privacy}
\label{section-computing-systems-local-ai}

One reason I care about hardware now is
local AI.

A local model runs on my own machine.
This can improve privacy because the model
does not inherently require sending the
prompt and documents to an external
server.

But a local model does not by itself prove
that the whole workflow is private.  The
surrounding application may still use
telemetry, cloud synchronization, web
search, or remote APIs.

So the privacy question should always be:

\begin{quote}
Which data leave my machine, through which
program, and for what reason?
\end{quote}

For local AI I should learn to distinguish:

\begin{itemize}
\item model size;
\item numerical precision and
quantization;
\item system RAM;
\item GPU memory or unified memory;
\item CPU inference and GPU inference;
\item local applications and cloud APIs.
\end{itemize}

A powerful laptop can run useful local
models, but a laptop is constrained by
weight, battery, heat, and power.  I should
not expect a portable machine to be a
permanent substitute for a future
workstation.

\section{The Apple question}
\label{section-computing-systems-apple}

My worry about Apple should also be stated
precisely.

The concern is not simply that macOS is
proprietary.  I am worried about becoming
dependent on a system in which the same
company controls the hardware, operating
system, account ecosystem, cloud services,
and an increasing number of AI features.

There are several separate questions.

\medskip\noindent
{\bf Can I install my own software?}
On macOS I can already install a great deal
of software outside Apple's App Store.
For mathematical and programming work this
includes most of the ordinary Unix-style
toolchain.

So saying that Apple simply prevents me
from installing software is too crude a
description of the problem.

\medskip\noindent
{\bf Can I replace the operating system?}
This depends on the particular Mac and on
the state of Linux support for its
hardware.  Even when Linux can run, not all
hardware features necessarily work equally
well.

Thus I should never buy a particular Mac
on the assumption that it will later
become an ideal Linux laptop without first
checking that exact model.

\medskip\noindent
{\bf Can I upgrade the machine?}
This is a much clearer limitation.
Modern Mac laptops integrate many
components very tightly.  RAM and storage
are generally not ordinary user-upgradable
modules.

That means I must buy enough memory and
storage at the beginning, and hardware
obsolescence can force replacement of more
of the machine than I would like.

\medskip\noindent
{\bf Can Apple force me to use its AI?}
No general principle forces my computing
workflow to use Apple's AI services.  I
can use third-party software and local
models.

The more serious strategic worry is
ecosystem dependence: if my files,
workflows, devices, authentication, and
services become tightly tied to one
vendor, leaving becomes expensive even if
no single restriction prevents me from
leaving.

\medskip\noindent
{\bf Privacy.}
Privacy is not the same as openness.
Closed software may have strong privacy
properties, and open-source software may
still send data to remote services if it
is configured to do so.

The robust habit is therefore not to trust
a brand-level slogan.  I should understand
where my data live, what is synchronized,
which services receive it, and which parts
of my workflow can work completely
offline.

\section{Mac versus Linux for a laptop}
\label{section-computing-systems-mac-linux}

The choice is not good hardware versus bad
hardware.

A Mac laptop may give an unusually strong
combination of battery life, performance,
screen, trackpad, speakers, compactness,
and quiet operation.  Those are real
advantages for a person whose entire
computer has to travel in a backpack.

A Linux laptop can give much more control
over the operating system and, depending
on the machine, much better repairability
and upgradeability.

Linux also gives direct experience with
the environment used by a huge amount of
scientific computing, servers, and AI
infrastructure.

Therefore the question for my next laptop
is not which platform is morally better.
It is which compromises matter during the
period in which this laptop has to be my
entire computing system.

\section{What I should compare before buying}
\label{section-computing-systems-buying}

For every serious candidate, I should make
the same checklist.

\begin{enumerate}
\item CPU performance for sustained work.
\item Amount of RAM and maximum RAM.
\item Whether RAM is replaceable.
\item SSD capacity and replaceability.
\item GPU and memory available to the GPU.
\item Quality of Linux support.
\item Battery life under real workloads.
\item Weight including the charger.
\item Cooling and fan noise.
\item Screen, keyboard, and trackpad.
\item Availability of replacement parts.
\item Ease of battery replacement.
\item Expected operating-system support.
\item Ability to run useful local AI.
\item Total price with enough RAM and
storage for several years.
\end{enumerate}

The right comparison is between complete
machines configured for my actual needs,
not between the cheapest advertised
versions.

\section{What I should learn next}
\label{section-computing-systems-next}

A sensible learning order is:

\begin{enumerate}
\item CPU, RAM, SSD, GPU, and VRAM.
\item Processes, filesystems, and memory.
\item The shell and the Unix filesystem.
\item Package managers.
\item Permissions and users.
\item Processes and services.
\item SSH and remote computing.
\item Partitions and filesystems.
\item Booting, UEFI, and firmware.
\item Drivers and the Linux kernel.
\item Backups and disk encryption.
\item Containers and virtual machines.
\item Local AI runtimes and model formats.
\end{enumerate}

I do not need to master all of this before
buying a computer.  The point is to stop
treating the machine as a mysterious
sealed object and gradually understand
each layer well enough to control my own
workflow.
